\documentclass[../../main.tex]{subfiles}

\begin{document}

\subsection{Il \textit{Far Detector}} \label{subsec:il-far-detector}
  \begin{figure}[ht]
    \centering
    \includegraphics[width=.5\textwidth]{DUNE/dune-project/fd}
    \caption{
      Illustrazione delle caverne sotterranee per il FD di DUNE e i sistemi
      criogenici; sono indicati l'asse e la direzione del fascio.
      \cite{abi2020deepundergroundneutrinoexperiment}
    }
    \label{fig:fd}
  \end{figure}
  Il \textit{Far Detector} è una camera a proiezione temporale ad argon liquido
  (LArTPC) modulare che combina una massa fiduciale nell'ordine dei $10^6$~kg,
  necessaria per la ricerca di eventi rari, con la capacità di ricostruire le
  immagini di tali eventi con una risoluzione spaziale millimetrica.
  Questa caratteristica permette di distinguere le firme dei processi fisici dal
  rumore di fondo.\\
  Un ulteriore punto di forza del rivelatore risiede nelle eccellenti proprietà
  di scintillazione dell'argon liquido, che emette una luce di scintillazione
  veloce a una lunghezza d'onda di 126,8~nm (UV).
  La luce di scintillazione, una volta spostata nello spettro visibile,
  viene raccolta dai rivelatori di fotoni (\textit{Photon Detectors}) (PD).
  Questi PD forniranno un tempo ``zero'' per ogni evento registrato dalla camera
  a proiezione temporale (TPC), indicando il momento in cui gli elettroni di
  ionizzazione iniziano la deriva (\textit{drift}).
  Confrontando questo tempo di inizio con l'istante in cui il segnale di
  ionizzazione raggiunge l'anodo, è possibile ricostruire la topologia
  dell'evento, inoltre, i PD forniranno informazioni calorimetriche.\\
  La struttura del FD è illustrata in figura~\ref{fig:fd}.


\end{document}